\documentclass[11pt,a4paper,oneside,openany]{book}
\usepackage{biblioflow-ista}

\renewcommand{\titolodocumento}{Test Incident Report}
\renewcommand{\sottotitolodocumento}{Test Incident Report}

% Intestazione di un incidente: componente, severità, data.
\newcommand{\datiincidente}[3]{%
  \textbf{Componente:} #1\par
  \textbf{Severità:} #2\par
  \textbf{Data:} #3\par\medskip}

\begin{document}
\frontmatter
\copertina
\tableofcontents

\mainmatter
\chapter{Introduzione}

Il presente documento costituisce l'Incident Report del progetto BiblioFlow.
In questo documento vengono analizzate le problematiche identificate durante
l'esecuzione della suite di test che caratterizza la versione di partenza e
durante la sua riesecuzione come suite di regressione. Nell'ambito di questa
analisi vengono pertanto presi in esame tutti i livelli di testing coinvolti.

\section{Obiettivi del Documento}

L'obiettivo primario è fornire una mappatura chiara delle criticità
individuate nella suite di test, distinguendo i difetti del sistema da quelli
dei test stessi. Questa attività di revisione e correzione è propedeutica al
miglioramento della qualità complessiva del sistema e all'aumento
dell'affidabilità della suite di test, garantendo una base solida per gli
sviluppi futuri.

\chapter{Incident Report}

In questa sezione vengono dettagliati gli incidenti specifici affrontati
durante la fase di validazione. Per ogni anomalia viene riportata la
descrizione dell'errore, l'analisi delle cause e l'intervento correttivo
applicato.

\section{Test Incident n.1}

\datiincidente{Test di Sistema -- esecuzione della suite pre-modifica sul tag
\texttt{baseline-pre-cr-bf-01}}{Minore (bloccante per l'esecuzione)}{28 agosto 2026}

\subsection{Descrizione del Problema}

La prima esecuzione della suite di caratterizzazione è stata tentata
direttamente sul tag che congela la versione di partenza. Il comando non è
neppure partito:

\begin{lstlisting}
$ git checkout baseline-pre-cr-bf-01
$ npm test
npm error Missing script: "test"
\end{lstlisting}

La versione di partenza non conteneva alcuna infrastruttura di test: niente
Vitest, niente script \texttt{test}, nessun database di test né fixture. Il
tag, che per definizione non può essere modificato, non era quindi
eseguibile come riferimento dei test.

\subsection{Risoluzione Tecnica}

Invece di modificare il tag, l'infrastruttura di test è stata aggiunta in un
commit dedicato (\texttt{7469969}) che non tocca i flussi da caratterizzare.
Il confronto con la baseline lo dimostra: nelle cartelle \file{src} e
\file{prisma} cambiano soltanto quattro file, e solo per correzioni
richieste dal linter (un tipo esplicito al posto di \texttt{any}, due
apostrofi trasformati in entità HTML e un commento per il linter).

\begin{lstlisting}
$ git diff --stat baseline-pre-cr-bf-01 7469969 -- src prisma
 src/app/api/admin/richieste/route.ts            | 8 +++++++-
 src/app/libri/[id]/page.tsx                     | 2 +-
 src/app/prenota/page.tsx                        | 2 +-
 src/components/accessibility/live-announcer.tsx | 5 +++++
 4 files changed, 14 insertions(+), 3 deletions(-)
\end{lstlisting}

Il commit aggiunge gli script \texttt{test} e
\texttt{test:db:prepare}, che crea un PostgreSQL effimero e rifiuta qualunque
indirizzo che non identifichi esplicitamente un database di test.

\subsection{Esito}

La suite è stata eseguita sul commit \texttt{7469969}: 5 file e 34 casi su 34
conformi all'oracolo. Nei documenti di test il riferimento di esecuzione è
indicato esplicitamente accanto al tag, così che chiunque possa riprodurre il
risultato.

\section{Test Incident n.2}

\datiincidente{Test di Integrazione --
\file{tests/pre-modifica/automation-notifications.test.ts}}{Minore
(ambiente di esecuzione)}{4 settembre 2026}

\subsection{Descrizione del Problema}

Durante la riesecuzione dei casi invariati come regressione, la suite delle
automazioni e delle notifiche non riusciva a preparare il database. Prima di
ogni esecuzione il setup sincronizza lo schema del database di test lanciando
la CLI di Prisma:

\begin{lstlisting}
const result = spawnSync(executable, ["db", "push"], {
  cwd: process.cwd(),
  env: { ...process.env, DATABASE_URL: databaseUrl },
  encoding: "utf8",
  shell: process.platform === "win32",
});
\end{lstlisting}

Su Windows il file eseguibile è \texttt{prisma.cmd} e viene lanciato
attraverso la shell. La cartella di lavoro conteneva uno spazio nel percorso:
la shell spezzava il percorso dell'eseguibile in due parole e il comando non
veniva trovato. Il fallimento non dipendeva né dai test né dal sistema.

\subsection{Risoluzione Tecnica}

Per non alterare né i test né l'applicazione durante una verifica di
regressione, la stessa copia di lavoro è stata esposta tramite un collegamento
(junction) a un percorso senza spazi e i cinque casi interessati sono stati
rieseguiti da lì, filtrandoli per identificativo:

\begin{lstlisting}
npx vitest run tests/pre-modifica/automation-notifications.test.ts \
  -t "PRE-AUT-00[124]|PRE-NOT-00[12]"
\end{lstlisting}

\subsection{Esito}

I cinque casi sono risultati conformi alla baseline. L'incidente è stato
classificato come problema dell'ambiente: la pipeline di integrazione
continua, che gira su Linux, non ne è affetta.

\section{Test Incident n.3}

\datiincidente{Test di Integrazione --
\file{tests/pre-modifica/automation-notifications.test.ts}}{Maggiore (falso
esito positivo)}{12 settembre 2026}

\subsection{Descrizione del Problema}

Durante la correzione del middleware, che da quel momento verifica davvero il
token di sessione tramite Auth.js, è emerso che l'intera suite delle
automazioni e delle notifiche non veniva più eseguita da tempo. Il modulo
\file{lib/auth.ts}, importato a cascata dalle route sotto test, carica
\texttt{next-auth}; questo importa \texttt{next/server}, e il resolver ESM di
Node non riesce a risolverlo perché il pacchetto di Next non dichiara il campo
\texttt{exports}:

\begin{lstlisting}
Error: Cannot find module '.../node_modules/next/server'
Did you mean to import "next/server.js"?
\end{lstlisting}

L'errore avveniva al caricamento del file: Vitest segnava i sei casi come
\emph{skipped}, non come falliti. Il riepilogo della suite restava verde e
nessuno se ne era accorto.

\subsection{Risoluzione Tecnica}

\textbf{Codice Errato:} la configurazione lasciava a Node il caricamento dei
pacchetti di Auth.js.

\begin{lstlisting}
test: {
  environment: "node",
  include: ["tests/**/*.test.ts", "tests/**/*.spec.ts"],
  passWithNoTests: true,
  ...
},
\end{lstlisting}

\textbf{Codice Corretto:} i due pacchetti passano dalla trasformazione di
Vite, che risolve correttamente l'estensione del modulo. Non cambia ciò che
viene testato, solo come viene caricato.

\begin{lstlisting}
test: {
  environment: "node",
  ...
  server: {
    deps: {
      inline: [/next-auth/, /@auth\/core/],
    },
  },
},
\end{lstlisting}

Una volta eseguita davvero, la suite ha rivelato che il caso PRE-NOT-002
chiamava l'elenco delle notifiche senza sessione: nel frattempo la route era
stata corretta di proposito per usare sempre l'utente autenticato invece di
un parametro \texttt{userId} preso dall'indirizzo. Al test è stato aggiunto il
mock della sessione, così che continui a verificare ciò che gli interessa:
filtro, paginazione e conteggi.

\subsection{Esito}

La suite torna a essere eseguita: sei casi su sei passano. L'incidente ha
mostrato che un riepilogo verde non basta: va controllato anche il numero di
casi effettivamente eseguiti, perché un file che non si carica produce casi
saltati e non falliti.

\section{Test Incident n.4}

\datiincidente{Test di Sistema -- \texttt{TC-PRE-016}
(\file{tests/pre-modifica/admin-ui.test.ts})}{Minore}{12 settembre 2026}

\subsection{Descrizione del Problema}

Il caso TC-PRE-016 era stato scritto per fotografare il comportamento della
baseline: con un cookie di sessione presente, il middleware lascia proseguire
la richiesta verso \texttt{/prenota}. Il test usava però un valore
\emph{inventato} (\texttt{token-baseline}) e si aspettava che passasse.

Quando il middleware è stato corretto per verificare firma e scadenza del
token, il test è fallito: la verifica del token è asincrona e un valore
inventato porta ora al login, mentre il test chiamava il middleware in modo
sincrono e pretendeva che la richiesta proseguisse.

Il fallimento non indicava una regressione. Il test non descriveva un
requisito ma il difetto stesso: in produzione bastava inviare un cookie
qualsiasi (\texttt{authjs.session-token=x}) per ottenere dati reali. Un test
scritto così avrebbe fatto apparire la correzione come una regressione.

\subsection{Risoluzione Tecnica}

\textbf{Codice Errato:} il test accettava un cookie inventato e verificava una
chiamata sincrona.

\begin{lstlisting}
it("[TC-PRE-016] lascia proseguire la UI prenotazione con cookie di sessione", () => {
  const response = appMiddleware.middleware(
    request("http://localhost/prenota", "GET", undefined,
      "authjs.session-token=token-baseline"));
  expect(response.status).toBe(200);
});
\end{lstlisting}

\textbf{Codice Corretto:} il caso è stato capovolto e reso asincrono, come la
verifica reale del token; un nuovo caso usa un token firmato davvero, perché
senza il caso \enquote{utente autenticato} anche un middleware che rifiuta
tutti sembrerebbe corretto.

\begin{lstlisting}
it("[TC-PRE-016] un cookie di sessione inventato non apre piu' la UI prenotazione", async () => {
  const response = await appMiddleware.middleware(
    request("http://localhost/prenota", "GET", undefined,
      "authjs.session-token=token-baseline"));
  expect(response.status).toBe(307);
  const destinazione = new URL(response.headers.get("location") as string);
  expect(destinazione.pathname).toBe("/login");
});
\end{lstlisting}

\subsection{Esito}

Il caso è classificato come \textbf{modificato} rispetto alla baseline, con
la motivazione documentata nel test stesso. La correzione ha confermato che i
test di caratterizzazione vanno riletti quando il difetto che fotografano
viene corretto: altrimenti diventano un vincolo che protegge il difetto.

\end{document}
